\subsection{A height bound at every arithmetic point}
\label{sec:actual-point-height}

We now use the corrected fibres on the actual model to bound the
prescribed regularized height, including at boundary points. All
degrees and heights are normalized by the degree of the field of
definition. Auxiliary field extensions are allowed, but no bound for
their degrees is assumed or needed.

Retain the original host $E$, $f=[E^+:\mathbf Q]$, signature $(p,q)$,
selected rank $m\in\{p,q\}$, and cutoff $B$. Write
\[
 \Lambda_B=\log B+\log\operatorname{rd}(E)
       +\left|L'(1,\chi_{-d})/L(1,\chi_{-d})\right|,
 \qquad X_B=2\sum_{\ell\in S_B}\frac{\log\ell}{\ell-1}
                         +2\log\operatorname{rd}(E).
 \tag{PH1}
\]
Here $S_B=S_0\cup S_{\rm new}(B)$ is the original full set of
correction primes. The set $S_{\rm new}(B)$ is the prescribed index
set for the additional vertical factors. We use
its proved bound $X_B\le16\Lambda_B$ and assume
$B\ge\max\{128,4f(p+q)^2\}$, as well as the original initial-cutoff
conditions, in particular $|S_0|\le B$ and
$S_B\subset S_0\cup\{\ell\text{ prime}:\ell\le B\}$.
Choose the original admitted normal model with its locally free
corrected modules first, using
Proposition~\ref{prop:actual-hodge-free-model}; form its prescribed
determinant coordinates and vertical factors on that model, and
retain them throughout the argument. This order is part of the
original construction. Flattening a previously nonfree module is
not being asserted to preserve its old determinant: it can introduce
an exceptional vertical divisor. If these modules were already
locally free, their later pullbacks are exact by that proposition.

\paragraph{The point and its determinant.}
Fix a point, extend its field $L$ if necessary, and write $G_z$ for
the original semiabelian fibre and $A_z$ for its abelian quotient.
Let $\Lambda$ be the integral tangent lattice, let $I$ be the
cocharacter lattice of its split generic torus, embedded in the
tangent space by differentiation at one, and let $r$ be its
$\mathcal O_E$-rank. Thus $0\le r\le m$. The integral
arithmetic-point construction gives an exact sequence onto its image
\[
 0\longrightarrow I\otimes_{\mathbf Z}\mathcal O_L
   \longrightarrow\Lambda\longrightarrow V\longrightarrow0,
 \qquad V\subset\operatorname{Lie}(\mathcal A_z/\mathcal O_L),
 \tag{PH2}
\]
where $\mathcal A_z$ is the semistable N\'eron identity component.
This is the original sequence. Over each normal valuation ring,
the closure of the generic torus is a closed smooth torus, and its
generic splitting extends uniquely \cite[Propositions 3.3.1.5
and 3.3.1.7]{point-height-lan}. Its tangent map is injective on
the special fibre, so its integral tangent image is saturated.
The quotient map extends to the N\'eron model by its mapping
property. Its generic tangent kernel is the torus tangent space.
Saturation therefore identifies its integral kernel with
$I\otimes\mathcal O_L$, proving \textup{(PH2)}. Passing to an
eigenlattice takes the image actually obtained, without assuming
that the image is saturated. These identifications are preserved
by semistable base change \cite{conrad-semistable,actual-point-height-source}.

Let $\Lambda_\tau$ be a selected integral eigenlattice and
$V_\tau$ its image in the quotient. Choose $v_1,\ldots,v_r\in I$
independent over $E$ with finite index
\[
                  i_T=[I:\textstyle\sum_a\mathcal O_Ev_a].
\]
Choose $q_{\tau,1},\ldots,q_{\tau,m-r}\in\Lambda_\tau$
whose images form an $L$-basis of the quotient eigenspace. Then
$b_\tau=(e_\tau v,q_\tau)$ is an $L$-basis of the full
eigenspace, where $e_\tau$ is the coefficient projector.
For a full lattice $M$ and an $L$-basis $b$, put
\[
 J(M,b)=\log[M:C]-\log[\mathcal O_L b:C],
 \tag{PH3}
\]
where $C$ is any common sublattice of finite index. This signed
quantity is independent of $C$. With the prescribed quotient
metric $\pi H$, set
\[
 \begin{split}
 d_\tau(F)&=\frac{1}{[L:\mathbf Q]}
   \left(J(F_\tau,b_\tau)
     -\frac12\sum_{\iota:L\hookrightarrow\mathbf C}
                     \log\det\operatorname{Gram}_\iota(q_\tau)\right),\\
 \Phi_B(z)&=\sum_\tau d_\tau(F)-\frac12\log i_T
                                  +\frac{fr}{2}\log(4\pi).
 \end{split}
 \tag{PH4}
\]
Here $F_\tau$ is the fibre of the original corrected bundle.
These formulas are the original labelled determinant and point
quantity, including the signs of the toric index and $4\pi$ term.

\begin{proposition}[The actual determinant bound]
\label{prop:actual-point-determinant-bound}
For every original arithmetic point, including the boundary,
\[
 \Phi_B(z)\le fmX_B+T-fr\log\operatorname{rd}(E)
                                  +\frac{fr}{2}\log(4\pi)
                    \le52fm\Lambda_B,
 \qquad T=\sum_\tau\widehat{\deg}_{\rm n}V_\tau.
 \tag{PH5}
\]
The same bound covers $r=m$ and zero-dimensional abelian quotients.
\end{proposition}

\begin{proof}
We first compute the toric determinant, rather than discarding its
finite index. The trace-dual basis description of the coefficient
projectors identifies the saturated $\tau$-eigenlattice of
$I\otimes\mathcal O_L$ with the $\tau$-image of the determinant
ideal of $I$, multiplied by $\tau(\mathfrak D_E)^r$ on determinants.
To see the factor, first take $I=\mathcal O_E$. In the decomposition
$E\otimes L=\prod_\sigma L$, the integral elements supported on
the $\tau$ factor are precisely
$\tau(\mathfrak D_E)\mathcal O_L e_\tau$. This follows by writing
$e_\tau$ in a trace-dual integral basis; the inverse different is
the trace dual of $\mathcal O_E$. Localization proves the same
identity for every projective $I$, and taking determinants proves
the rank-$r$ assertion. Consequently
\[
 \frac{J((I\otimes\mathcal O_L)_\tau,e_\tau v)}{[L:\mathbf Q]}
       =-r\log\operatorname{rd}(E)+\frac{\log i_T}{2f}.
 \tag{PH6}
\]
The formula also holds for $r=0$, with $i_T=1$.

The exact sequence of lattices in \textup{(PH2)}, restricted to the
chosen eigenlattice and its actual image, gives an exact sequence
of determinant lines. Equivalently, choose local integral bases
for the kernel and image and lift the latter: the change-of-basis
matrix is block triangular. At infinity the quotient block has
the stated Gram determinant. Thus
\[
 d_\tau(\Lambda)
   =-r\log\operatorname{rd}(E)+\frac{\log i_T}{2f}
                              +\widehat{\deg}_{\rm n}V_\tau.
 \tag{PH7}
\]
Theorem~\ref{thm:actual-hodge-specialization} gives
$\mathfrak a_\tau F_\tau\subset\Lambda_\tau\subset F_\tau$
and $\log N\mathfrak a_\tau/[L:\mathbf Q]\le X_B$.
Locally a quotient of a rank-$m$ lattice killed by this ideal
has length at most $m$ times its valuation. Hence
\[
       0\le d_\tau(F)-d_\tau(\Lambda)\le mX_B.
\]
Sum \textup{(PH7)} over the $f$ labels. The term $\log i_T/2$
cancels the negative term in \textup{(PH4)}, proving the first
inequality of \textup{(PH5)}.

Corollary~\ref{cor:actual-selected-point-degrees} gives
$T\le36f(m-r)\Lambda_B$, with both integral indices retained.
Also $X_B\le16\Lambda_B$, $\log\operatorname{rd}(E)\ge0$,
and $\log(4\pi)\le\Lambda_B$. Therefore the right side of
\textup{(PH5)} is at most
$f\{16m+36(m-r)+r/2\}\Lambda_B\le52fm\Lambda_B$.
No slope estimate is applied when the quotient rank is zero.
\end{proof}

\paragraph{The boundary metric on the same determinant.}
We give the limiting calculation needed to compare \textup{(PH4)}
with the regularized metric. On an original normal-crossing chart,
choose the symplectic period coordinates of the semiabelian family.
Their period map has the form
\[
 M(z)(u+\tau(z)v),\qquad
 \tau(z)=T_0(z)+\sum_{a=1}^{h}\frac{\log z_a}{2\pi i}B_a,
 \tag{PH8}
\]
where $M(0)$ is invertible, $M,T_0$ are holomorphic at the centre,
and the $B_a$ are symmetric integral positive semidefinite matrices.
The polarization fixes the symplectic form. The extension of the
semiabelian family identifies the bounded periods with its original
central periods and the radical with its actual torus. The positive
definite form on the abelian quotient is the one from the same
polarization. To recall their construction, the central periods continue
holomorphically and span the central tangent space. Their rational
period subspace is coisotropic. Choose a Lagrangian in it fixed by
all monodromies and extend it to a symplectic basis. The resulting
monodromies have the upper triangular form with integral symmetric
blocks $B_a$. Positivity gives positive semidefiniteness of those
blocks; subtracting their logarithms gives the holomorphic $T_0$
by the removable-singularity argument for these weight-one periods.
Central spanning makes $M(0)$ invertible. Thus these are the period
coordinates of the original boundary construction
\cite{actual-point-height-source}.

Here is the exact cancellation that makes the limit valid for all
rates of approach. Put $K=\ker\sum_aB_a$ and let $\Pi$ be the real
orthogonal projector onto $K$. Positive semidefiniteness gives
$B_a\Pi=\Pi B_a=0$ for every $a$. In period coordinates the raw
Hermitian form is $Y(z)^{-1}$, where $Y=\operatorname{Im}\tau$.
Its quotient by the toric space $K^\perp$ is represented by
\[
 C(z)=\Pi\bigl(\Pi\operatorname{Im}T_0(z)\Pi+1-\Pi\bigr)^{-1}\Pi.
 \tag{PH9}
\]
Indeed $\Pi Y\Pi=\Pi\operatorname{Im}T_0\Pi$. For a vector $y$,
the vector $p=YC(z)y$ satisfies $p-y\in K^\perp$ and is
$Y^{-1}$-orthogonal to $K^\perp$; its pairing with $y'$ is
$y'^*C(z)y$. This proves the formula without an asymptotic error.
The restriction of $\operatorname{Im}T_0(0)$ to $K$ is positive
definite, since it is the Riemann form of the abelian quotient.
Thus \textup{(PH9)} and its transport by $M(z)$ extend
continuously to the centre, including $K=0$ and $K=\mathbf R^g$.

Let $s_\tau(z)$ be the actual selected tangent frames. Use the exact
moving toric period vectors, project them to the coefficient
eigenspaces, and complete them by continuous transverse vectors.
The coefficient projectors vary continuously, so this is an
eigenframe whose change-of-basis determinant from $s_\tau(z)$
is continuous and nonzero at the centre. The determinant of its
Gram matrix factors exactly as
\[
 \det\operatorname{Gram}(\hbox{toric vectors})\,
           \det\operatorname{Gram}_{C(z)}(\hbox{transverse vectors}).
 \tag{PH10}
\]
The coefficient eigenspaces are orthogonal. Thus the quotient by
the toric part in one eigenspace is the restriction of the full
toric quotient. This is the determinant formula for an orthogonal projection,
or the Schur-complement identity. If $n_z(I)$ denotes the original
product of toric Gram determinants divided by $i_T$, division by
$n_z(I)$ cancels the first factors in \textup{(PH10)} exactly.
The remaining factors converge by \textup{(PH9)} and the
nonzero frame determinants. No estimate of a small holomorphic
error times a divergent logarithm is being made.

At the centre write $b_\tau=C_\tau s_\tau$. A cocharacter
derivative has period $2\pi i$ times that derivative. Converting
the raw quotient form to the prescribed metric $\pi H$ therefore
gives the following limiting squared determinant norm:
\[
 \nu_\iota=i_T\prod_\tau
       \frac{(4\pi)^{-r}\det\operatorname{Gram}_\iota(q_\tau)}
                                {|\iota(\det C_\tau)|^2}>0.
 \tag{PH11}
\]
In detail, each selected label contributes $\pi^m$ from the
full tangent metric, $\pi^{-(m-r)}$ from the quotient Gram
matrix, and $|2\pi i|^{-2r}$ from the toric period vectors.
Their product is $(4\pi)^{-r}$. This verifies the normalization
in \textup{(PH11)} directly.

The original cusp correction $U$ satisfies
$U(z)\ge-\log n_z(I)$ for every saturated isotropic $I$.
For the zero lattice this says $U\ge0$. It is the original
first-minimum bound, with the same counting function.
The corrected interior logarithmic weight is the Hodge weight
plus $U/2$. Formulae \textup{(PH9)--(PH11)} show that its lower
limit at the point is at least the algebraic frame weight plus
$\tfrac12\log\nu_\iota$. Call this latter value $x_\iota$.
Changing projective coordinate charts only adds the prescribed
continuous algebraic frame term, so the statement holds in the
original homogeneous coordinates.

\begin{theorem}[The original all-point height floor]
\label{thm:actual-all-point-height}
Let $\overline{\mathcal L}_{B,a/b,t}$ be the prescribed smooth
regularization with positive rational perturbation $a/b$ by the
original ample line $\overline{\mathcal A}$. There is a finite
$K_{\mathcal A}\ge0$, independent of the point, such that for
every fixed $a/b>0$ and all sufficiently small $t>0$,
\[
 h_{\overline{\mathcal L}_{B,a/b,t}}(P)
        \ge-52fm\Lambda_B-(a/b)K_{\mathcal A}
                  \qquad(P\in X(\overline{\mathbf Q})).
 \tag{PH12}
\]
The threshold is chosen before $P$ and its field of definition.
All finite vertical factors and the original complex metrics
are retained. Thus \textup{(PH12)} supplies the all-point
hypothesis of Proposition~\ref{geo-reciprocal-height-limit}.
\end{theorem}

\begin{proof}
First verify the arithmetic metric identity. The determinant
tensor of the inverse corrected bundles is the original
coordinate line of degree $\mu_B$, with its prescribed vertical
correction factors. This can be checked on its local frames:
if $T_{\tau,v}$ changes a rational frame to an integral frame
of $F_\tau$ at a finite place, the transition of the tensor is
the product of $\det T_{\tau,v}^{-1}$. The original coordinate
and vertical transitions have exactly this product. Their ratio
on an arithmetic point is one element $g\in L^\times$, obtained
from the original algebraic frame factor multiplied by its
vertical factors. In particular this is a single algebraic
element at every embedding and every finite place.

More explicitly, choose nonzero homogeneous coordinates $w_{k,c_k}$
at the point and put
\[
 R_v=\prod_k\left(\frac{\max_j|w_{k,j}|_v}{|w_{k,c_k}|_v}
                         \right)^{\mu_{B,k}}.
\]
For the original finite factors write
$t_\ell=(v_{{\rm num},\ell}/v_{{\rm den},\ell})
\prod_k w_{k,c_k}^{-v_{\ell,k}}$ at the point. These are
nonzero: the normalized rational ratio is a unit at every
characteristic-zero point of the original chart. With its
algebraic Hodge frame factor $u_L$, which already includes the
initial $S_0$ correction, the original local
determinant presentation gives exactly
\[
 \left(\prod_\tau|\det T_{\tau,v}|_v\right)R_v=|g|_v,
 \qquad g=u_L\prod_{\ell\in S_{\rm new}(B)}t_\ell.
 \tag{PH13a}
\]
At infinity its vertical factor is
$|\iota(t_\ell)|\prod_k|\iota(w_{k,c_k})|^{v_{\ell,k}}$,
and $\mu_{B,k}=a_{\mu,k}+\sum_{\ell\in S_{\rm new}(B)}v_{\ell,k}$.
These identities verify the transition assertion for the exact
coordinate line, including every finite factor.

The signed index formula is
\[
 J(F_\tau,b_\tau)
   =\sum_{v\nmid\infty}\log|\det T_{\tau,v}|_v
       +\sum_{\iota:L\hookrightarrow\mathbf C}
                                      \log|\iota\det C_\tau|.
 \tag{PH13}
\]
Normalize the finite absolute values so the product formula
sums over all complex embeddings, counted separately.
The finite transition identity just described and
\textup{(PH11)} give
\[
                    h_B^0(P)+\frac1{[L:\mathbf Q]}
                         \sum_\iota x_\iota=-\Phi_B(P).
 \tag{PH14}
\]
Here $h_B^0$ consists of the finite homogeneous-coordinate height
of degree $\mu_B$ and the archimedean logarithms of the
additional vertical factors indexed by $S_{\rm new}(B)$. To check the signs, substitute
\textup{(PH13)} and
\[
 \tfrac12\log\nu_\iota
  =\tfrac12\log i_T-\tfrac{fr}{2}\log(4\pi)
    +\tfrac12\sum_\tau\log\det\operatorname{Gram}_\iota(q_\tau)
    -\sum_\tau\log|\iota\det C_\tau|.
\]
The terms in $C_\tau$ cancel; the remaining algebraic scalar
$g$ cancels by the product formula. What remains is precisely
the negative of \textup{(PH4)}. This proves
\textup{(PH14)} on the actual determinant line.

Let $\varphi$ be the original bounded upper semicontinuous
potential, relative to the fixed smooth reference metric. The
interior is dense. The preceding boundary lower-limit comparison
and upper semicontinuity give
\[
                 x_\iota\le\log h_{\rm ref}(w_\iota)
                                      +\varphi(P_\iota).
 \tag{PH15}
\]
The prescribed smooth regularizations satisfy
$\varphi\le\varphi_t$ at every point for all sufficiently small
$t>0$. For clarity, their threshold is global: a finite cover by
the original inner charts has a positive common margin;
submean inequalities for the circle-invariant smoothing kernel
and the prescribed positive correction give domination on each
inner chart. Their maximum construction gives domination on the
whole compact space. The same margin and smoothness conditions
are fixed before evaluating any point
\cite{actual-point-height-source}.

The exact degree and metric formula for the prescribed rational
line now reads
\[
 h_{\overline{\mathcal L}_{B,a/b,t}}(P)
   =h_B^0(P)+\frac1{[L:\mathbf Q]}\sum_\iota
        \{\log h_{\rm ref}(w_\iota)+\varphi_t(P_\iota)\}
                                      +(a/b)h_{\overline{\mathcal A}}(P).
 \tag{PH16}
\]
The ample height is bounded below: compare a positive tensor
power of $\mathcal A$ with its projective-coordinate height,
whose reference height is nonnegative; the logarithm of the
ratio of continuous positive metrics is bounded on the compact
complex space. This supplies $K_{\mathcal A}$, which may
depend on the fixed model and metric but not on $P$.
Combining \textup{(PH14)--(PH16)} with \textup{(PH5)} proves
\textup{(PH12)}. Properness realizes every algebraic point on the
model, after the harmless field extensions already allowed.
Normalized base change preserves all terms. If $m=0$, the
selected determinant is the unit line and the assertion follows
from the ample-height bound alone.
\end{proof}

The all-point floor is now an established input. The signed
intersection limit, the original weighted mean-height comparison,
and the finite rank losses must still be assembled on the same
family. No whole geometric proof-map node is closed by this
pointwise result alone.
